% Appendix A — Formal definitions and proofs
\section{Formal definitions}
\label{app:formal}

\paragraph{Innovation game} An innovation game is a tuple
$\mathcal G=(\mathcal S,\mathcal L,\mathcal V,u,\mathcal I,\beta)$: a state space $\mathcal S$ of
encoded problems; a typed solution language $\mathcal L$ whose well-typed members are complete
candidates; a verifier $\mathcal V:\mathcal S\times\mathcal L\to\mathcal Y\cup\{\bot\}$ returning an
outcome or a failure; a utility $u:\mathcal Y\cup\{\bot\}\to[-1,1]$ with $u(\bot)=-1$; an incumbent
map $\mathcal I:\mathcal S\to\mathcal L$ giving the best verified solution known for a state before
play; and a declared budget $\beta$ (candidates, evaluations, simulations, wall time). A family
$F$ is a distribution over $\mathcal S$ induced by one source or generator; families sharing a source
form a dependency component. A \emph{one-shot game} ($H=1$) ends after one move. A \emph{sequential
game} of horizon $T$ appends each verified outcome to the state,
$s_{t+1}=s_t\oplus(a_t,\mathcal V(s_t,a_t))$.

\paragraph{The four conditions} The innovation-game hypothesis applies to a problem class when
(i) every instance can be encoded as $s\in\mathcal S$ without the verifier's private data;
(ii) $\mathcal L$ is decidable---well-typedness is checked without executing the candidate;
(iii) $\mathcal V$ is independent of the player: it runs outside the player's process, from inputs the
player cannot write; and (iv) $\beta$ is declared before play.

\paragraph{Information sets and the Zero contract} Let $\Omega(s)$ be the information available to the
player at commitment and $c=\mathrm{hash}(a)$ the commitment. The Zero contract
(\cref{def:zero}) holds when the supervisor's counters satisfy
$K=|\{a:\text{committed}\}|=1$, $E=|\{\text{candidate-specific evaluations before } c\}|=0$ and
$R=|\{\text{revisions after any feedback}\}|=0$, and $\Omega(s)$ contains no function of
$\mathcal V(s,\cdot)$. Track C0 has $\Omega(s)=s$; track CH has
$\Omega(s)=s\oplus h$ for a frozen history $h$ identical across arms; track SEQ applies the contract at
each step; track S relaxes $K$, $E$ and $R$ to a declared budget.

\paragraph{Utility and estimand} For a state $s$ with reference $b(s)$ and domain loss $e$,
$U=\mathrm{clip}(\log(b(s)/e),-1,1)$ for valid moves and $U=-1$ for failures. For arm $m$ and
comparator $c$, the family-level paired difference is
$\delta_F=\mathbb E_{s\sim F}[U_m(s)-U_c(s)]$ and the estimand is
$\Delta=\mathbb E_{F}[\delta_F]$ over confirmation families with equal family weights.

\paragraph{Search-equivalent budget} Let $\mu_{\mathrm{search}}(B)$ be the family-weighted mean
utility of the declared search procedure at budget $B$ on a measured grid $\mathcal B$. Then
$B_{\mathrm{eq}}=\min\{B\in\mathcal B:\mu_{\mathrm{search}}(B)\ge\mu_{\mathrm{EZ}}(1)\}$; if no
$B\in\mathcal B$ qualifies, we report $B_{\mathrm{eq}}>\max\mathcal B$ and do not extrapolate. Its
interval is the set of grid points selected across bootstrap replicates (\cref{app:statistics}).

\subsection{Proofs}

\paragraph{\Cref{prop:closed} (closed support)} Let $p_\theta(\cdot\mid s)$ have support
$\mathcal A_{\mathrm{train}}$ for every $\theta$ in the parameter set $\Theta$ reachable by training.
Then $\Pr_\theta[a\notin\mathcal A_{\mathrm{train}}]=\sum_{a\notin\mathcal A_{\mathrm{train}}}
p_\theta(a\mid s)=0$. Every update maps $\Theta$ into itself, so the support never grows.
\hfill\tikz\draw[line width=0.5pt] (0,0) rectangle (0.62em,0.62em);

\paragraph{\Cref{prop:repr} (representation limit)} Let $z=E(s)$. For any policy
$\pi(a\mid z)$, $\mathbb E[U\mid s,\pi]\le\sup_a\mathbb E[U\mid s,a]$ for every $s$, with equality only
if $\pi$ puts its mass on moves optimal for $s$. If $E(s)=E(s')$ and no move optimal for $s$ is optimal
for $s'$, $\pi(\cdot\mid z)$ cannot put all its mass on optimal moves for both, so it is suboptimal on
at least one. Taking expectations over states gives the inequality.
\hfill\tikz\draw[line width=0.5pt] (0,0) rectangle (0.62em,0.62em);

\paragraph{Decomposition} For any proposer $\pi$ and any larger language
$\mathcal L^\star\supseteq\mathcal L$, the gap to the best achievable move splits into three
non-negative terms:
\begin{align*}
\mathbb E\sup_{a\in\mathcal L^\star}\mathbb E[U\mid s,a]-\mathbb E[U\mid z,\pi]
={}&\underbrace{\mathbb E\sup_{a\in\mathcal L^\star}\mathbb E[U\mid s,a]-\mathbb E\sup_{a\in\mathcal L}\mathbb E[U\mid s,a]}_{\text{language}}\\
&+\underbrace{\mathbb E\sup_{a\in\mathcal L}\mathbb E[U\mid s,a]-\mathbb E\sup_{a\in\mathcal L}\mathbb E[U\mid z,a]}_{\text{encoder}}\\
&+\underbrace{\mathbb E\sup_{a\in\mathcal L}\mathbb E[U\mid z,a]-\mathbb E[U\mid z,\pi]}_{\text{policy}}.
\end{align*}
The
experiments estimate the encoder term with oracle-state controls, the language term with a richer
decoder on development data, and the policy term with search at large budget.
